\chapter{Software Requirement Specifications}
\label{ch:srs}

This chapter specifies the behavior required of the Manzil research prototype. Requirements are testable where current evidence permits; numerical values that require benchmarking are explicitly identified as design targets. The specification avoids account management, enforcement, and other modules not justified by the approved scope.

\section{System Features}

The prototype is divided into six feature groups:

\begin{enumerate}
\item \textbf{Scenario management}: select a network, demand profile, controller, seed, incident, duration, and emergency mission.
\item \textbf{Traffic perception}: process licensed recorded video into detections, tracks, lane-region observations, and aggregate measures with timestamps and confidence.
\item \textbf{Simulation and normal control}: step SUMO through a single owner, read traffic state, apply fixed or adaptive actions, and preserve legal phase transitions.
\item \textbf{Emergency routing and corridor}: validate origin and destination, compute dynamic paths, apply stable rerouting, activate moving priority, detect passage, and recover.
\item \textbf{Monitoring}: show map state, signals, traffic indicators, ambulance route, corridor state, warnings, and data-source mode.
\item \textbf{Experiment records}: persist immutable configuration, run events, aggregate metrics, software versions, seed, and termination state for comparison and export.
\end{enumerate}

\section{Functional Requirements}

\begin{longtable}{p{0.10\textwidth}p{0.78\textwidth}}
\caption{Functional requirements}\label{tab:functional-requirements}\\
\toprule
\textbf{ID} & \textbf{Requirement}\\
\midrule\endfirsthead
\toprule
\textbf{ID} & \textbf{Requirement}\\
\midrule\endhead
FR-01 & The system shall validate a scenario configuration before starting and shall reject missing network, controller, seed, duration, or demand inputs.\\
FR-02 & The system shall start, pause, resume, and stop one controlled SUMO run while preserving an unambiguous run state.\\
FR-03 & The simulation adapter shall be the sole owner of SUMO stepping and signal actuation.\\
FR-04 & The perception worker shall process a selected recorded video and emit timestamped detections, tracks, source metadata, and processing status.\\
FR-05 & The observation service shall aggregate valid vehicle observations into defined road, lane, or approach intervals without silently mixing sources.\\
FR-06 & The operator shall be able to select fixed-time, rule-adaptive, or the configured RL controller before a run.\\
FR-07 & Every signal action shall identify its intersection, requested and applied phase, effective simulation step, controller, and reason.\\
FR-08 & The system shall enforce minimum green and configured yellow/all-red clearance before incompatible movements receive right of way.\\
FR-09 & The emergency service shall validate ambulance origin and destination and return a clear no-route response if either is invalid or disconnected.\\
FR-10 & The route service shall compute a valid remaining path using nonnegative time-based edge costs and shall record the cost snapshot.\\
FR-11 & The rerouting service shall re-evaluate the remaining path after a closure, invalidated edge, or sustained cost improvement above the configured threshold.\\
FR-12 & The corridor coordinator shall activate priority only for a bounded look-ahead set of intersections on the current ambulance route.\\
FR-13 & The corridor coordinator shall detect passage or cancellation and restore the affected intersection to a safe normal-control state.\\
FR-14 & A route revision shall cancel priority requests that no longer belong to the remaining path before activating new ones.\\
FR-15 & The dashboard shall display the current source mode, simulation status, controller, signal state, traffic indicators, ambulance position, route, and corridor state.\\
FR-16 & The dashboard shall present stale, invalid, disconnected, and fallback states visibly rather than continuing with apparently current data.\\
FR-17 & The system shall record configuration, seed, component versions, periodic metrics, route revisions, signal events, failures, and termination reason for each run.\\
FR-18 & The comparison service shall align selected runs by metric definition and warn when network, demand, duration, warm-up, or measurement versions are incompatible.\\
FR-19 & The operator shall be able to export run configuration and metrics in a documented machine-readable format.\\
FR-20 & On loss of the learned controller or invalid action, the control service shall fall back to a configured safe baseline and record the transition.\\
\bottomrule
\end{longtable}

\section{Quality Attributes and Non-Functional Requirements}

\begin{longtable}{p{0.11\textwidth}p{0.22\textwidth}p{0.57\textwidth}}
\caption{Non-functional requirements and design targets}\label{tab:nfr}\\
\toprule
\textbf{ID} & \textbf{Attribute} & \textbf{Requirement or target}\\
\midrule\endfirsthead
\toprule
\textbf{ID} & \textbf{Attribute} & \textbf{Requirement or target}\\
\midrule\endhead
NFR-01 & Reproducibility & A run shall store network and demand identifiers, seed, controller configuration, software/model versions, warm-up, duration, and metric schema.\\
NFR-02 & Safety in simulation & Incompatible green movements shall never be commanded intentionally; all emergency transitions shall use the configured clearance sequence.\\
NFR-03 & Auditability & Applied decisions shall be distinguishable from requests, and their simulation timestamps and reasons shall be retained.\\
NFR-04 & Freshness & Each observation shall carry produced-at and effective simulation time. Stale thresholds are configuration values and stale data shall not silently drive control.\\
NFR-05 & Responsiveness & The control action for step $t$ should be available before its configured decision deadline. The exact deadline is a design target to be set after baseline profiling, not an achieved value.\\
NFR-06 & Isolation & Background video inference and dashboard clients shall not advance SUMO or directly actuate signals.\\
NFR-07 & Usability & A trained operator should be able to configure the core demonstration from one experiment screen, with units and validation messages visible.\\
NFR-08 & Portability & Services should be runnable from versioned environments; containerization is proposed after local baseline integration is stable.\\
NFR-09 & Maintainability & Component interfaces shall use typed, versioned data contracts and controller implementations shall share one interface.\\
NFR-10 & Data integrity & A completed run configuration shall be immutable; derived summaries shall retain their metric version and source run identifiers.\\
NFR-11 & Privacy & Video shall be limited to licensed or team-authorized material, stored only as needed, and excluded from exports unless explicitly requested.\\
NFR-12 & Availability & The prototype shall fail visibly and recover to a safe stopped or baseline-controlled simulation state; no public-service availability is claimed.\\
\bottomrule
\end{longtable}

\section{Assumptions and Dependencies}

SUMO and TraCI are required for controlled traffic experiments \cite{lopez2018sumo,sumodocs}. A selected network has legal lane connections and signal plans after import and cleanup. Video sources have sufficient quality and permission for the selected evaluation. The simulation uses one clock owner. Edge costs are nonnegative. Emergency origin and destination lie on ambulance-accessible edges. The dashboard is an operator tool, not an autonomous public dispatch system.

React and FastAPI are fixed by the approved proposal for the interface and application service boundary \cite{reactdocs,fastapidocs}. Python is implied by the planned perception, RL, and TraCI ecosystem. PostgreSQL, a WebSocket transport, Docker, and specific chart or map libraries are proposed dependencies and remain replaceable until implementation confirms them.

\section{Hardware and Software Requirements}

\begin{table}[H]
\centering
\caption{Development and demonstration environment}
\label{tab:environment}
\begin{tabularx}{\textwidth}{p{0.26\textwidth}YY}
\toprule
\textbf{Area} & \textbf{Development/training need} & \textbf{Demonstration need}\\
\midrule
Compute & Multi-core CPU; CUDA-capable GPU is desirable for detector/RL experiments but exact memory is unconfirmed & CPU able to run selected SUMO network and services; GPU only if live video inference is demonstrated\\
Memory/storage & Sufficient for environments, video, checkpoints, run logs, and repeated traces; capacity to be measured & Selected scenario, service builds, model, and demonstration records\\
Core software & Python, SUMO/TraCI, detector/tracker framework, RL library, Node.js/React toolchain, Git & Version-pinned runtime and browser\\
Data tools & OpenStreetMap import, SUMO NetEdit, annotation tooling & Prepared network and approved demonstration clips\\
Persistence & File-backed prototype acceptable initially; PostgreSQL proposed for structured records & Same schema/version as evaluation records\\
\bottomrule
\end{tabularx}
\end{table}

\section{Use Cases}

The primary actor is a research operator. SUMO, the video source, and the experiment store are supporting systems. Evaluators are read-only consumers of dashboards and exported records. Figure \ref{fig:use-case} presents the coverage; Tables \ref{tab:uc1}--\ref{tab:uc8} define the principal flows.

\begin{figure}[H]
\centering
\begin{tikzpicture}[font=\footnotesize,uc/.style={draw,ellipse,minimum width=3.5cm,minimum height=0.8cm,align=center},actor/.style={draw,rounded corners,fill=manzilblue!10,minimum width=1.8cm,minimum height=0.8cm},node distance=5mm and 11mm]
\node[actor] (op) {Research operator};
\node[uc,right=20mm of op,yshift=18mm] (cfg) {Configure scenario};
\node[uc,below=of cfg] (run) {Run/control simulation};
\node[uc,below=of run] (video) {Process traffic video};
\node[uc,below=of video] (inspect) {Inspect traffic state};
\node[uc,right=17mm of cfg] (dispatch) {Dispatch ambulance};
\node[uc,below=of dispatch] (corr) {Monitor route/corridor};
\node[uc,below=of corr] (reroute) {Handle route change};
\node[uc,below=of reroute] (compare) {Compare/export runs};
\foreach \x in {cfg,run,video,inspect}{\draw (op)--(\x);}
\draw (op.east) -- ++(5mm,0) |- (dispatch.west);
\draw (op.east) -- ++(5mm,0) |- (corr.west);
\draw (op.east) -- ++(5mm,0) |- (reroute.west);
\draw (op.east) -- ++(5mm,0) |- (compare.west);
\end{tikzpicture}
\caption{Proposed operator use-case diagram}
\label{fig:use-case}
\end{figure}

\newcommand{\uctable}[7]{%
\begin{table}[H]\centering
\caption{#2}\label{#1}
\begin{tabularx}{\textwidth}{p{0.22\textwidth}Y}
\toprule
\textbf{Name / actor} & #2 / Research operator\\
\textbf{Preconditions} & #3\\
\textbf{Trigger} & #4\\
\textbf{Main success flow} & #5\\
\textbf{Alternative/failure flow} & #6\\
\textbf{Postconditions} & #7\\
\bottomrule
\end{tabularx}\end{table}}

\uctable{tab:uc1}{UC-01 Configure a scenario}{Network and controller descriptors are available.}{Operator requests a new experiment.}{Select network, demand, controller, seed, warm-up, duration, source mode, incident, and optional mission; validate; store a draft configuration.}{Reject incompatible controller/network, invalid duration, missing seed, or inaccessible input with field-specific messages.}{A valid, versioned configuration is ready; no run has started.}

\uctable{tab:uc2}{UC-02 Run and control a simulation}{UC-01 is valid; SUMO adapter is available.}{Operator selects Start.}{Create run record; launch SUMO; complete warm-up; step the simulation; collect state; request and apply lawful actions; publish snapshots; stop at terminal condition.}{If launch/TraCI fails, mark failed and retain logs. Pause/resume does not create competing clock owners. Invalid controller action invokes fallback.}{Run is completed, stopped, or failed with an explicit reason.}

\uctable{tab:uc3}{UC-03 Process traffic video}{Licensed video and model configuration are available.}{Operator selects a video job.}{Validate source; decode frames; detect vehicles; associate tracks; map valid tracks to regions; aggregate interval observations; store metrics and job status.}{Unreadable media, unsupported format, model failure, missing calibration, or low-confidence observation is reported; physical queue length is omitted when calibration is unavailable.}{Completed observations identify source, model, interval, units, and confidence.}

\uctable{tab:uc4}{UC-04 Inspect traffic state}{A run or perception job has emitted data.}{Operator opens overview or intersection detail.}{Display source mode, freshness, counts, queues/proxies, speeds, phases, controller, and warnings; allow time/intersection selection.}{Disconnected or stale streams remain visibly marked; last value is not presented as current.}{No system decision is changed; the inspected state is traceable to its source.}

\uctable{tab:uc5}{UC-05 Dispatch an ambulance}{A running emergency-enabled scenario exists.}{Operator submits origin and destination.}{Validate accessible edges; create mission; compute route/cost; assign simulated ambulance; publish route; request bounded corridor look-ahead.}{Invalid edge or no route returns a reason without partial priority. Duplicate active mission is rejected in the core single-ambulance mode.}{Mission is active with a valid current route, or rejected cleanly.}

\uctable{tab:uc6}{UC-06 Monitor route and corridor}{An emergency mission is active.}{Ambulance advances through the network.}{Update position and remaining path; request safe transition ahead; display corridor states; detect passage; release and recover each intersection.}{Missing vehicle, stale position, unsafe request, or corridor timeout triggers cancellation/fallback and a visible alert.}{Mission progress and all priority transitions are recorded.}

\uctable{tab:uc7}{UC-07 Handle blockage or changed route}{An active mission has a feasible current position.}{Closure or sustained cost change triggers evaluation.}{Recompute path from a valid remaining edge; compare benefits; if threshold is met, atomically revise route, cancel obsolete requests, and activate new look-ahead.}{If improvement is too small, retain route. If no route exists, cancel upcoming priority and mark mission blocked.}{Exactly one route version is current; every revision has cause and cost snapshot.}

\uctable{tab:uc8}{UC-08 Compare and export runs}{At least two completed or failed run records exist.}{Operator selects runs and metrics.}{Check comparability; align metric definitions; present distributions and configuration differences; export selected records.}{Warn or refuse direct comparison when network, demand, duration, warm-up, or metric version differs. Failed runs are not silently omitted.}{Comparison and export identify source run IDs and schema version.}

\section{Proposed Graphical User Interface}
\label{sec:gui}

The following figures are proposed wireframes, not screenshots of implemented software. Their purpose is to validate information hierarchy and required operator actions without fabricating progress.

\begin{figure}[H]
\centering
\begin{tikzpicture}[x=1cm,y=1cm,font=\scriptsize]
\draw[rounded corners] (0,0) rectangle (14,7.2);
\fill[manzilblue] (0,6.5) rectangle (14,7.2); \node[text=white,anchor=west] at (.3,6.85) {MANZIL / Traffic Overview};
\draw[fill=manzilgray] (.3,.4) rectangle (9.3,6.2); \node at (4.8,5.9) {Network map: congestion, signals, ambulance and route};
\foreach \x/\y/\c in {1.4/4.8/manzilgreen,3.4/3.7/manzilamber,5.5/4.5/manzilred,7.6/2.6/manzilgreen}{\fill[\c] (\x,\y) circle (.16);}
\draw[very thick,manzilblue] (1.4,4.8)--(3.4,3.7)--(5.5,4.5)--(7.6,2.6);
\draw (9.6,4.7) rectangle (13.7,6.2); \node[align=left] at (11.65,5.45) {Run RUN-021\\Controller: Rule adaptive\\Source: SUMO / current};
\draw (9.6,2.7) rectangle (13.7,4.5); \node[align=left] at (11.65,3.6) {Emergency EV-01\\ETA: computed\\Corridor: TRANSITION\\Next: I-07};
\draw (9.6,.4) rectangle (13.7,2.5); \node[align=left] at (11.65,1.45) {Warnings / event timeline\\stale and fallback states\\remain visible};
\end{tikzpicture}
\caption{Proposed traffic overview wireframe}
\label{fig:wire-overview}
\end{figure}

\begin{figure}[H]
\centering
\begin{tikzpicture}[x=1cm,y=1cm,font=\scriptsize]
\draw[rounded corners] (0,0) rectangle (14,6.7);
\fill[manzilblue] (0,6) rectangle (14,6.7); \node[text=white,anchor=west] at (.3,6.35) {Intersection I-07 / Perception Detail};
\draw[fill=black!10] (.3,2.1) rectangle (8.4,5.7); \node at (4.35,5.35) {Processed video preview (when available)};
\draw[manzilgreen,thick] (1.2,2.8) rectangle (2.3,4.0); \node at (1.75,2.55) {track 17};
\draw[manzilamber,thick] (3.1,3.0) rectangle (4.4,4.2); \node at (3.75,2.75) {track 23};
\draw (8.7,4.1) rectangle (13.7,5.7); \node[align=left] at (11.2,4.9) {Source: video-03\\Model: version / checksum\\Freshness: 120 ms\\Calibration: unavailable};
\draw (8.7,2.1) rectangle (13.7,3.8); \node[align=left] at (11.2,2.95) {North: count 12; queue proxy 7\\East: count 8; queue proxy 3\\Confidence and invalid flags};
\draw (.3,.3) rectangle (13.7,1.8); \node at (7,1.05) {Time-series area: count, occupancy, queue proxy, processing throughput};
\end{tikzpicture}
\caption{Proposed perception and intersection-detail wireframe}
\label{fig:wire-perception}
\end{figure}

\begin{figure}[H]
\centering
\begin{tikzpicture}[x=1cm,y=1cm,font=\scriptsize]
\draw[rounded corners] (0,0) rectangle (14,6.3);
\fill[manzilblue] (0,5.6) rectangle (14,6.3); \node[text=white,anchor=west] at (.3,5.95) {Emergency Dispatch and Route};
\draw (.3,3.9) rectangle (5.0,5.3); \node[align=left] at (2.65,4.6) {Origin [select edge]\\Destination [select edge]\\Dispatch [button]};
\draw[fill=manzilgray] (5.3,.4) rectangle (13.7,5.3); \node at (9.5,5.0) {Map with route alternatives and current mission};
\draw[very thick,manzilgreen] (6.2,1.4)--(7.6,2.5)--(9.3,2.0)--(11.2,3.4)--(12.7,4.1);
\draw[thick,dashed,manzilamber] (6.2,1.4)--(8.5,3.6)--(10.4,4.2)--(12.7,4.1);
\draw (.3,1.9) rectangle (5,3.6); \node[align=left] at (2.65,2.75) {Selected cost: seconds\\Route version: 3\\Reroute threshold: config\\Reason: sustained improvement};
\draw (.3,.4) rectangle (5,1.6); \node[align=left] at (2.65,1.0) {Corridor lifecycle\\I-07 ACTIVE / I-08 REQUESTED};
\end{tikzpicture}
\caption{Proposed emergency dispatch and route wireframe}
\label{fig:wire-dispatch}
\end{figure}

\begin{figure}[H]
\centering
\begin{tikzpicture}[x=1cm,y=1cm,font=\scriptsize]
\draw[rounded corners] (0,0) rectangle (14,6);
\fill[manzilblue] (0,5.3) rectangle (14,6); \node[text=white,anchor=west] at (.3,5.65) {Experiment Configuration};
\draw (.5,4.35) rectangle (8.4,4.85); \node[anchor=west] at (.7,4.6) {Network and demand profile};
\draw (.5,3.70) rectangle (8.4,4.20); \node[anchor=west] at (.7,3.95) {Controller and policy version};
\draw (.5,3.05) rectangle (8.4,3.55); \node[anchor=west] at (.7,3.30) {Seed, warm-up and duration};
\draw (.5,2.40) rectangle (8.4,2.90); \node[anchor=west] at (.7,2.65) {Incident and emergency mission};
\draw (.5,1.75) rectangle (8.4,2.25); \node[anchor=west] at (.7,2.00) {Metric schema and sample interval};
\draw (8.8,2.0) rectangle (13.5,4.85); \node[align=left] at (11.15,3.4) {Validation summary\\\checkmark IDs resolve\\\checkmark phase plan legal\\! video calibration absent\\! RL checkpoint mismatch};
\draw[fill=manzilgreen!25] (9.0,.7) rectangle (11.0,1.4); \node at (10,1.05) {Save draft};
\draw[fill=manzilblue!20] (11.3,.7) rectangle (13.3,1.4); \node at (12.3,1.05) {Start run};
\end{tikzpicture}
\caption{Proposed experiment-configuration wireframe}
\label{fig:wire-config}
\end{figure}

\begin{figure}[H]
\centering
\begin{tikzpicture}[x=1cm,y=1cm,font=\scriptsize]
\draw[rounded corners] (0,0) rectangle (14,6);
\fill[manzilblue] (0,5.3) rectangle (14,6); \node[text=white,anchor=west] at (.3,5.65) {Comparative Analytics};
\draw (.4,4.15) rectangle (13.6,5.0); \node at (7,4.58) {Run selector / comparability warning / metric definition};
\draw (.4,.4) rectangle (8.7,3.85); \node at (4.55,3.55) {Distribution with mean and variability};
\draw[fill=manzilblue!35] (1.3,.8) rectangle (2.4,2.8); \draw[fill=manzilgreen!35] (3.2,.8) rectangle (4.3,2.2); \draw[fill=manzilamber!50] (5.1,.8) rectangle (6.2,1.8);
\node at (1.85,.6) {Fixed}; \node at (3.75,.6) {Rule}; \node at (5.65,.6) {RL};
\draw (9.0,.4) rectangle (13.6,3.85); \node[align=left] at (11.3,2.1) {Configuration differences\\Seed coverage\\Failed runs\\Export data + manifest};
\end{tikzpicture}
\caption{Proposed comparative-analytics wireframe; bars are schematic and not results}
\label{fig:wire-analytics}
\end{figure}

\section{Data Design}
\label{sec:data-design}

Persistent data is reserved for configuration, observations, decisions, missions, and outcomes that require audit or comparison. The live road graph, current TraCI subscription, and frame buffers are in-memory runtime state. Raw video remains an external file with a content identifier rather than a database blob.

\begin{figure}[H]
\centering
\begin{tikzpicture}[font=\scriptsize,node distance=8mm and 9mm,ent/.style={draw,rounded corners,fill=manzilgray,minimum width=3.1cm,minimum height=1.0cm,align=left},rel/.style={-{Latex},thick}]
\node[ent] (scenario) {SCENARIO\\id, network, demand};
\node[ent,right=of scenario] (run) {RUN\\id, seed, status, versions};
\node[ent,right=of run] (obs) {OBSERVATION\\source, time, measures};
\node[ent,below=of run] (decision) {CONTROL DECISION\\requested/applied, reason};
\node[ent,below=of scenario] (mission) {MISSION\\origin, destination, status};
\node[ent,right=of mission] (route) {ROUTE VERSION\\edges, cost, trigger};
\node[ent,right=of route] (priority) {PRIORITY EVENT\\intersection, state, time};
\node[ent,below=of decision] (metric) {METRIC SAMPLE\\name, value, unit, version};
\draw[rel] (scenario)--node[above]{1:N}(run);
\draw[rel] (run)--node[above]{1:N}(obs);
\draw[rel] (run)--(decision); \draw[rel] (run)--(metric);
\draw[rel] (run)--(mission); \draw[rel] (mission)--node[above]{1:N}(route);
\draw[rel] (route)--(priority);
\end{tikzpicture}
\caption{Proposed experiment-record entity relationship model}
\label{fig:er}
\end{figure}

\begin{longtable}{p{0.19\textwidth}p{0.22\textwidth}p{0.48\textwidth}}
\caption{Data dictionary}\label{tab:data-dictionary}\\
\toprule
\textbf{Entity} & \textbf{Key fields} & \textbf{Integrity rule}\\
\midrule\endfirsthead
\toprule
\textbf{Entity} & \textbf{Key fields} & \textbf{Integrity rule}\\
\midrule\endhead
Scenario & scenario\_id, network\_version, demand\_version & References immutable network and demand manifests.\\
Run & run\_id, scenario\_id, seed, controller\_version, status & Configuration becomes immutable after start; terminal reason required.\\
Observation & run\_id/job\_id, source\_mode, entity\_id, observed\_at, valid\_until, values, units & Source and units are mandatory; validity interval cannot be inferred.\\
ControlDecision & run\_id, intersection\_id, requested\_phase, applied\_phase, effective\_step, reason & Requested and applied actions are separate fields.\\
Mission & mission\_id, run\_id, vehicle\_id, origin, destination, status & Only one active core mission per run.\\
RouteVersion & mission\_id, version, edge\_list, cost\_seconds, trigger, created\_at & Version increases monotonically; path begins at valid current position.\\
PriorityEvent & mission\_id, route\_version, intersection\_id, state, requested/applied time & References the route version that justified priority.\\
MetricSample & run\_id, metric\_name, metric\_version, value, unit, interval & Comparisons require matching metric name, version, and unit.\\
Artifact & run/job\_id, type, URI, checksum, license/provenance & Content is addressable and provenance is retained.\\
\bottomrule
\end{longtable}

\section{Risk Analysis}

\small
\begin{longtable}{p{0.20\textwidth}p{0.10\textwidth}p{0.10\textwidth}p{0.28\textwidth}p{0.13\textwidth}}
\caption{Risk register}\label{tab:risks}\\
\toprule
\textbf{Risk} & \textbf{Likelihood} & \textbf{Impact} & \textbf{Mitigation} & \textbf{Fallback}\\
\midrule\endfirsthead
Map import lacks valid signals/turns & Medium & High & Inspect and clean a small extract early; validate routes and phases & Use a documented synthetic grid\\
Local video lacks labels or permission & Medium & High & Use licensed public sequences and record provenance & Evaluate public data only; state domain gap\\
Occlusion causes unstable counts & High & Medium & Region rules, temporal aggregation, tracker comparison & Limit claims to labelled conditions\\
RL fails to converge or generalize & Medium & High & Stable baselines, bounded action space, tune only on training scenarios & Report baseline comparison; treat RL as negative/partial finding\\
Reward improves one metric unfairly & Medium & High & Multi-metric evaluation and starvation constraints & Reject policy violating safety or fairness checks\\
Rerouting oscillates & Medium & Medium & Benefit threshold, persistence, cooldown, closure exception & Freeze current feasible route\\
Preemption leaves unsafe phase & Low & Very high & Central state machine and invariant tests & Cancel request and use safe baseline transition\\
Service timing drifts from simulation & Medium & High & One step owner, simulation timestamps, bounded queues & Run synchronously without live perception\\
Insufficient compute & Medium & Medium & Profile early; smaller network/model and offline training & Demonstrate recorded inference and compact scenarios\\
Results are not reproducible & Medium & High & Immutable manifests, seeds, versioned metrics and exports & Exclude untraceable runs\\
\bottomrule
\end{longtable}
\normalsize

\section{Requirements Traceability}

\begin{longtable}{p{0.15\textwidth}p{0.20\textwidth}p{0.27\textwidth}p{0.20\textwidth}}
\caption{Requirements traceability matrix}\label{tab:traceability}\\
\toprule
\textbf{Requirements} & \textbf{Use cases} & \textbf{Design modules} & \textbf{Planned verification}\\
\midrule\endfirsthead
\toprule
\textbf{Requirements} & \textbf{Use cases} & \textbf{Design modules} & \textbf{Planned verification}\\
\midrule\endhead
FR-01, 17--19; NFR-01,10 & UC-01, UC-08 & Experiment service, store & Schema, validation, replay and comparison tests\\
FR-02,03,07; NFR-02,06 & UC-02 & Simulation adapter, control orchestrator & Single-owner and phase-invariant tests\\
FR-04,05; NFR-04,11 & UC-03, UC-04 & Perception worker, state aggregator & Annotated sequence evaluation and stale-data tests\\
FR-06,08,20 & UC-02 & Controller interface, safety filter & Fixed scenarios, illegal-action and fallback tests\\
FR-09--11,14 & UC-05, UC-07 & Graph/routing service & Disconnected graph, closure, hysteresis and path-validity tests\\
FR-12,13 & UC-06, UC-07 & Corridor coordinator & Lifecycle, timeout, cancellation and recovery tests\\
FR-15,16 & UC-04, UC-06 & API gateway, dashboard & UI state and warning acceptance checks\\
NFR-03,05,07--09,12 & All & Cross-cutting contracts and deployment & Log audit, profiling, usability walkthrough, failure injection\\
\bottomrule
\end{longtable}

\section{Conclusion}

The specification centers on traceability and failure visibility. It requires one simulation-clock owner, legal signal transitions, stable route revision, explicit source freshness, reproducible run records, and comparisons that cannot hide incompatible configurations. These requirements become procedural and architectural choices in the next two chapters.
